\section{Methodology}
\label{sec:methodology}

Our analysis proceeds through a six-stage pipeline: task and dataset
definition, model training, representation extraction, intrinsic dimension
estimation, simplicial complex construction, and persistent homology
computation. 

% ------------------------------------------------------------------
\subsection{Task Formulation and Datasets}
\label{subsec:dataset}
% ------------------------------------------------------------------

\paragraph{Task formulation.}
We study two modular arithmetic tasks over $\mathbb{Z}_p$ with prime
modulus $p = 97$, formulated as next-token prediction following
Power et al.~\cite{GrokkingOpenAI}. Each input is the four-token
sequence $
    \langle\, a \;\circ\; b \;=\; \rangle$ ,
where $a, b \in \{0, \ldots, p-1\}$ and $\circ \in \{+, \times\}$.
The model's vocabulary has size $|\mathcal{V}| = p + 2 = 99$,
comprising the $p$ residues, one operator token, and one equality token.
The training objective is to predict, at the position immediately
following the equality token, the single token corresponding to
\begin{equation}
    y = (a \;\circ\; b) \bmod p \;\in\; \{0, \ldots, 96\}.
    \label{eq:target}
\end{equation}
The loss is the cross-entropy between the model's output distribution
over $\mathcal{V}$ and the one-hot encoding of $y$; accuracy is the
fraction of examples for which the model's argmax equals $y$.

\paragraph{Modular addition.}
The full dataset consists of all $p^2 = 9{,}409$ ordered pairs $(a,b)$.
For example, the input $\langle\,45\;+\;72\;=\;\rangle$ should elicit
$y = (45 + 72) \bmod 97 = 20$. The task defines an abelian group on
$\mathbb{Z}_{97}$, admitting an exact decomposition via discrete Fourier
analysis over the group~\cite{nanda2023progressmeasuresgrokkingmechanistic}, which the model must
discover implicitly through gradient descent.

\paragraph{Modular multiplication.}
The dataset again contains all $p^2 = 9{,}409$ ordered pairs. For
example, $\langle\,15\;\times\;9\;=\;\rangle$ should elicit
$y = (15 \times 9) \bmod 97 = 38$. The non-zero residues
$\mathbb{Z}_{97}^*$ form a cyclic group of order $p-1 = 96$ under
multiplication; the case $a = 0$ or $b = 0$ trivially yields $y = 0$.
Multiplication is empirically harder to grok than addition, requiring
more training epochs before generalization occurs~\cite{GrokkingOpenAI}.

\paragraph{Train/validation splits.}
For each task and run, the $p^2$ examples are randomly partitioned into
a training set of fraction $q \in \{10\%, 15\%, 20\%, 25\%, 30\%\}$
and a held-out validation set of the remaining $1-q$ fraction~\footnote{No data augmentation or curriculum
ordering is applied.}. Varying $q$ controls the difficulty of
generalization: smaller fractions provide less coverage of the input
space, making it harder for the model to discover the underlying
algebraic structure. 

% ------------------------------------------------------------------
\subsection{Model Architecture and Training}
\label{subsec:model}
% ------------------------------------------------------------------

We use the decoder-only transformer of Power et al.~\cite{GrokkingOpenAI}
without modification: two transformer blocks, each with 4-head
self-attention and a feedforward sublayer of hidden dimension 512,
preceded by a 128-dimensional token-plus-positional embedding layer and
followed by a linear projection to $\mathbb{R}^{|\mathcal{V}|}$.
%
Training uses AdamW~\cite{loshchilov2018decoupled} with learning rate
$10^{-3}$, and weight decay
$\lambda = 1.0$, with full-batch gradient descent for $10^6 + 1$ epochs.
Validation accuracy is evaluated at each epoch without gradient updates.

% ------------------------------------------------------------------
\subsection{Representation Extraction}
\label{subsec:extraction}
% ------------------------------------------------------------------

At regular checkpoint epochs $t$, we perform a forward pass over the
\emph{entire} dataset with gradients disabled and record activations
at the position of the equality token from four components:
(1) the token-plus-positional embedding layer,
$A^{(0)}_t \in \mathbb{R}^{n \times 128}$;
(2) the first decoder block output,
$A^{(1)}_t \in \mathbb{R}^{n \times 128}$;
(3) the second decoder block output,
$A^{(2)}_t \in \mathbb{R}^{n \times 128}$; and
(4) the pre-softmax linear projection,
$A^{(3)}_t \in \mathbb{R}^{n \times 99}$,
where $n = p^2 = 9{,}409$. Each matrix $A^{(\ell)}_t$ is treated as a
point cloud in its ambient space, with one point per input example.
The equality-token position is used throughout because it is the
position from which the output distribution is computed, and is
therefore the most direct representation of the model's belief about
the answer at checkpoint $t$.

% ------------------------------------------------------------------
\subsection{TDA Pipeline}
\label{subsec:tda_pipeline}
% ------------------------------------------------------------------

\paragraph{Intrinsic dimension estimation.}
For each layer $\ell$ and checkpoint $t$, we estimate the intrinsic
dimension $\hat{d}^{(\ell)}_t$ of $A^{(\ell)}_t$ using the MLE of
Levina and Bickel~\cite{LevinaBickel} (Equation~\ref{eq:mle_id},
$k = 10$ nearest neighbors). Note that, this estimate serves two objectives: as a
standalone geometric descriptor tracked over training, and as the target dimension for the
subsequent reduction step.

\paragraph{Dimensionality reduction.}
We apply UMAP~\cite{McInnes2018} to reduce $A^{(\ell)}_t$ to $\tilde{A}^{(\ell)}_t \in \mathbb{R}^{n \times 2\hat{d}}$, with target dimension $2\hat{d}$ motivated by the Strong Whitney Embedding Theorem (Section~\ref{subsec:id}). As discussed in Section~\ref{subsec:id}, UMAP is a heuristic that preserves local connectivity but provides no formal topological guarantees; the reported Betti numbers are exact invariants of the constructed complex but should be interpreted as indicative of structural trends with respect to the underlying activation manifold.

\paragraph{Simplicial complex construction.}
On the reduced $\tilde{A}^{(\ell)}_t$, we construct a graph using the algorithm of Liu et al.~\cite{QuickMapper}: an edge $(x_i, x_j)$ is included if $\|x_i - x_j\|_2 \leq \varepsilon$, where $\varepsilon$ is the 95th percentile of distances to the $k=31$-th nearest neighbor across all points. The graph is expanded into a full simplicial complex via GUDHI's simplex tree~\cite{GUDHI}, adding 2- and 3-simplices wherever the corresponding cliques exist. The choices $k = 31$ and the 95th-percentile threshold are heuristic~\footnote{A sensitivity analysis over these choices is needed to establish the robustness of the reported signatures, and it remains future work}.

\paragraph{Persistent homology.}
We compute persistent homology via GUDHI's boundary matrix
reduction~\cite{GUDHI} for dimensions $k \in \{0, 1, 2\}$.
From each persistence diagram $\mathcal{D}^{(\ell,k)}_t$ we extract
$\beta^{(\ell,k)}_t$ (see Section~\ref{subsec:simplicial}) as the count of pairs with persistence exceeding
$5\%$ of the maximum observed distance (noise threshold $\tau = 0.05$).
We restrict to $k \leq 2$ because $\mathbb{Z}_{97} \cong$ the 97th
roots of unity embeds naturally in a 2-torus $\mathbb{T}^2$, whose
homology is trivial (i.e., equal to zero) above dimension 2; higher-dimensional features are
negligible in our experiments.

% ------------------------------------------------------------------
\subsection{Evaluation Protocol}
\label{subsec:eval}
% ------------------------------------------------------------------

We compare the longitudinal topological time series
$\{\beta^{(\ell,k)}_t,\, \hat{d}^{(\ell)}_t\}$ across all checkpoints
$t$, layers $\ell \in \{0,1,2,3\}$, and dimensions $k \in \{0,1,2\}$,
for each task and training fraction $q$. %Conditions with $q \geq 20\%$ (generalizing) are compared against $q \leq 15\%$ (non-generalizing).
We evaluate three hypotheses:

\begin{enumerate}[label=\textbf{(H\arabic*)}]
    \item \textbf{Betti stabilization.} In generalizing conditions,
          $\beta^{(\ell,0)}_t$ converges to a stable value as validation accuracy increases.

    \item \textbf{Layer-specific directionality.} Upon generalization,
          $\beta^{(\ell,0)}_t$ decreases in the decoder layers and
          increases in the linear layer.

    \item \textbf{Intrinsic dimension reduction.} Generalizing conditions
          exhibit a lower converged $\hat{d}^{(\ell)}_t$ across all
          layers than non-generalizing conditions.
\end{enumerate}

% Evaluation is conducted through visual analysis of longitudinal plots and tabular comparison of end-of-training values. The absence of formal statistical tests is acknowledged as Limitation (Section~\ref{subsec:conclusion_limitations}).